We wszystkich \ContentionRuns{} rundach testu (4 bazy, 13 strategii, 4 poziomy
współbieżności, 3 rozmiary, 3 powtórzenia) liczba rezerwacji w~bazie zgadzała
się z~liczbą zajętych miejsc. Liczba niespójnych rund wynosi
\ContentionInconsistent. Żadna strategia nie dopuściła do podwójnej sprzedaży,
więc wyniki różnią się wydajnością, a~nie poprawnością.

\paragraph{Valkey.} Skrypt Lua przy 10 mln biletów obsługuje
\cval{valkey}{10000000}{lua}{1}\,rezerwacji/s z~jednym klientem
i~\cval{valkey}{10000000}{lua}{1000}\,rez./s z~tysiącem klientów. Przepustowość
rośnie z~liczbą klientów, bo jednowątkowy serwer wykonuje skrypty jeden po
drugim bez blokad, a~sterownik valkey-go łączy równoległe wywołania w~potoki.
95.~percentyl czasu próby przy 1000 klientach to \cpn{valkey}{10000000}{lua}{1000}\,ms.
Strategia \texttt{watch} (optymistyczna, na całym zbiorze wolnych miejsc)
załamuje się: \ckn{valkey}{10000000}{watch}{1000} przerwanych transakcji na
rundę i~\cval{valkey}{10000000}{watch}{1000}\,rez./s. Ten sam silnik daje więc
wyniki różniące się o~kilka rzędów wielkości w~zależności od sposobu blokowania.

\paragraph{PostgreSQL.} Warunkowy \texttt{UPDATE} (\texttt{cond}) osiąga
maksimum przy 100 klientach (\cval{postgres}{10000000}{cond}{100}\,rez./s),
a~przy 1000 klientach spada do \cval{postgres}{10000000}{cond}{1000}\,rez./s
z~p95 równym \cpn{postgres}{10000000}{cond}{1000}\,ms. Tysiąc procesów
serwera na 10 rdzeniach czeka w~kolejce na blokady tych samych 500 wierszy.
\texttt{SELECT ... FOR UPDATE} przy 1000 klientach daje więcej
(\cval{postgres}{10000000}{forupdate}{1000}\,rez./s), choć przy 100 klientach
kolejność jest odwrotna. Przy dużej współbieżności sposób oczekiwania na
blokadę wiersza ma więc znaczenie.
\texttt{serializable} traci najwięcej na ponowieniach
(\ckn{postgres}{10000000}{serializable}{1000} konfliktów na rundę).

\paragraph{MariaDB.} Do 100 klientów MariaDB dorównuje PostgreSQL, a~z~\texttt{nowait}
osiąga nawet \cval{mariadb}{10000000}{nowait}{100}\,rez./s. Przy 1000 klientach
strategie czekające na blokadę spadają do kilkudziesięciu rezerwacji na
sekundę (\texttt{cond}: \cval{mariadb}{10000000}{cond}{1000}\,rez./s,
p95 \cpn{mariadb}{10000000}{cond}{1000}\,ms). Rundy trwają wtedy dziesiątki
sekund. Obserwujemy kolejki tysiąca wątków czekających na blokady rekordów
InnoDB w~trybie REPEATABLE READ. \texttt{nowait} nie czeka, tylko od razu
zgłasza konflikt, i~zachowuje \cval{mariadb}{10000000}{nowait}{1000}\,rez./s.
Przy wysokiej współbieżności lepiej szybko odrzucić próbę niż kolejkować.

\paragraph{MongoDB.} Z~jednym klientem MongoDB jest najwolniejszy
(\cval{mongo}{10000000}{atomic}{1}\,rez./s), bo każdy zapis czeka na dziennik.
Przepustowość rośnie z~liczbą klientów, bo zapisy wielu klientów dzielą
jeden zapis dziennika (\emph{group commit}). Przy 1000 klientach osiąga
\cval{mongo}{10000000}{atomic}{1000}\,rez./s, więcej niż PostgreSQL
i~MariaDB w~strategiach domyślnych. Transakcje wielodokumentowe (\texttt{txn})
są wolniejsze od pojedynczej atomowej aktualizacji dokumentu.

\paragraph{Rozmiar zbioru.} Wykresy dla 500 tys. i~10 mln biletów mają ten
sam kształt. Test dotyczy 500 gorących wierszy, więc decyduje o~nim sposób
blokowania, a~nie rozmiar tabeli.
